%!TEX root = ../thesis.tex
\chapter*{概要}
\thispagestyle{empty}
\begin{center}
  \scalebox{1.5}{基盤モデルを用いた}\\
  \scalebox{1.5}{シナリオ駆動型ナビゲーションの構築と評価}
\end{center}
\vspace{1.0zh}

本研究グループでは，人間が記述したシナリオとカメラ画像に基づき，移動ロボットを目的地まで移動させるシナリオ駆動型ナビゲーションを提案してきた．
この手法は走行時に自己位置推定を必要としないが，通路形状の認識と経路追従に用いる学習モデルは，走行対象の環境で収集したデータで学習する必要があった．
本研究では，通路分類と経路追従に基盤モデル（DINOv2，UniDepth，NoMaD）を応用し，走行対象の環境でデータ収集や再学習を行わずに，シナリオに従って目的地まで移動できるかを検証した．
基盤モデルの特徴抽出器を凍結し，追加した部分を主に学習することで，事前学習で獲得された表現を保つようにシステムを構成した．
シミュレータ上で，ある棟のデータのみでファインチューニングしたシステムを走行させた結果，その棟の7例に加え，データを収集していない別の棟の6例を含む13例すべてのシナリオで目的地に到達した．

\vspace{1.0zh}
% TODO: キーワードを決める．
\noindent キーワード：

\newpage
\chapter*{Abstract}
\thispagestyle{empty}
\begin{center}
  \scalebox{1.3}{Construction and Evaluation of Scenario-Driven}\\
  \scalebox{1.3}{Navigation Using Foundation Models}
\end{center}
\vspace{1.0zh}

We have proposed scenario-driven navigation, in which a mobile robot follows a human-written scenario with camera images, without localization.
This paper applies foundation models (DINOv2, UniDepth, and NoMaD) to its corridor classification and path following, and examines whether the robot can reach its destination without collecting data or retraining in the target environment.
In simulation, the robot, fine-tuned only on data from one wing of a building, reached the destination in all 13 scenarios, including six in another wing where no data were collected.

\vspace{1.0zh}
% TODO: Keywords
\noindent Keywords:
